\clearpage\begingroup
\SourceLeaf{133b}\HistoricalHead{128}{L}{PHYSIK.\quad NACHLASS.}
\begin{center}
\textbf{ERLÄUTERUNGEN ZU DEN ABSCHNITTEN [I.] BIS [III.]\textsuperscript{*)}.}\par\medskip
\textls[120]{Zum Abschnitt [I.], S. 117.}
\end{center}
\begingroup\renewcommand{\thefootnote}{}\footnotetext{*) Diese Erläuterungen beruhen im Wesentlichen auf Mitteilungen von \textsc{M. v. Rohr} und \textsc{H. Boegehold}.}\endgroup
\noindent\hspace*{\parindent}Der »\textit{Neue Algorithmus}« in [I.] ist im Wesentlichen der \textsc{Euler}sche, den \textsc{Gauss} auch in seinen \textit{Dioptrischen Untersuchungen} (Werke V, S. 250) benutzt hat. Die Formelgruppen [III] bis [VI] ergeben sich an der Hand der Figuren A und B durch Anwendung der gewöhnlichen Linsenformel:
\[
\frac1g+\frac1b=f
\]
\noindent $\left(g=\text{Gegenstandsweite},\ b=\text{Bildweite},\ \frac1f=\text{Brennweite}\right)$, indem man der Reihe nach eine Linse nach der andern durchrechnet. In Figur A wird der Gang der sogenannten »Hauptstrahlen«, in Figur B derjenige der sogenannten »Öffnungsstrahlen« durch das zentrierte System hindurch verfolgt.

Aus Figur B ergibt sich zunächst, da die Strahlen parallel einfallend gedacht werden, daß $\frac\alpha\beta=\frac1f$, also mit $\alpha=1$, daß $\beta=f$ ist, was die erste Gleichung [IV] liefert. Dann ist $\frac{\alpha'}\beta$ offenbar gleich $h-\frac1f$, wenn $h$ den Abstand zwischen der ersten und zweiten Linse bezeichnet. Daraus ergibt sich $\alpha'=hf-1$, in Übereinstimmung mit der zweiten Gleichung [III]; ferner sieht man, daß die Gegenstandsweite für die zweite Linse den Wert $h-\frac1f$ hat. Nach der Linsenformel folgt daraus die Beziehung:
\[
\frac1{h-\frac1f}+\frac{\beta'}{\alpha'}=f',
\]
\clearpage\SourceLeaf{134}\HistoricalHead{129}{R}{DIOPTRIK.}
\noindent oder nach Einsetzen der Werte:
\[
\frac f{fh-1}+\frac{\beta'}{fh-1}=f';
\]
\noindent dies liefert für $\beta'$
\[
\beta'=f'(fh-1)-f,
\]
\noindent was in der Tat mit der zweiten der Gleichungen [IV] übereinstimmt, u.s.w.

Ebenso ist es mit Figur A. Da $\delta$ gleich 1 gesetzt ist, so ist
\[
\frac{\gamma'}\delta=\gamma'=h,
\]
\noindent was mit der zweiten der Gleichungen [V] übereinstimmt; betrachtet man ferner in Figur A die Strecke $h$ als Gegenstandsweite für die zweite Linse, $\frac{\gamma'}{\delta'}$ als Bildweite, so liefert die Linsenformel wieder:
\[
\frac1h+\frac{\delta'}{\gamma'}=\frac1h+\frac{\delta'}h=f',
\]
\noindent oder:
\[
\delta'=hf'-1=(h,f'),
\]
\noindent was mit der zweiten der Gleichungen [VI] übereinstimmt. In dieser Weise verifiziert man leicht die gesamten Formelgruppen [III] bis [VI]. Auch findet man leicht folgende Rekursionsformeln, die für die späteren Rechnungen nützlich sind:
\[
\begin{array}{llll}
\alpha=1,&\beta=f,&\gamma=0,&\delta=1,\\
\alpha'=h\,\beta-\alpha,&\beta'=f'\,\alpha'-\beta,&\gamma'=h\,\delta-\gamma,&\delta'=f'\,\gamma'-\delta,\\
\alpha''=h'\beta'-\alpha',&\beta''=f''\,\alpha''-\beta',&\gamma''=h'\delta'-\gamma',&\delta''=f''\gamma''-\delta',\\
\alpha'''=h''\beta''-\alpha'',&\beta'''=f'''\alpha'''-\beta'',&\gamma'''=h''\delta''-\gamma'',&\delta'''=f'''\gamma'''-\delta'',
\end{array}
\]
\begin{center}u.s.w.\end{center}
\noindent In derselben Weise ergeben sich auch die Formeln [VII] bis [VIIc] für den farbigen Rand.
\begin{center}\textls[120]{Zum Abschnitt [II.], S. 120.}\end{center}
\noindent\hspace*{\parindent}Im Abschnitt [II.] sind drei Fernrohre mit Hilfe der in dem Abschnitt [I.] angegebenen Formeln analysiert, d.\,h. die Größen $\alpha,\beta,\gamma,\delta$ aus den gegebenen Stücken $h$ und $f$ bestimmt, daraus der farbige Rand. Um die Beziehung auf Abschnitt [I.] zu erleichtern, sind im \textsc{Gauss}schen Texte vielfach in eckigen Klammern [\ ] die Bezeichnungen von [I.] hinzugefügt, sodaß die Rechnung ohne Mühe verfolgt werden kann.

\newpage
Auf der Göttinger Sternwarte befindet sich ein \textsc{Fraunhofer}sches Zugfernrohr, dessen Objektivdeckel von \textsc{Gauss}' Hand den Vermerk trägt: »\textsc{Gauss} 1815«. Die Dimensionen desselben sind denen, die \textsc{Gauss} in Abschnitt [II, B] angibt, sehr ähnlich, aber doch nicht mit ihnen identisch. Eine sorgfältige Messung ergab z.\,B. für die Brennweite des Objektivs den Wert $437{,}1\,\mathrm{mm}$, während \textsc{Gauss} $444\,\mathrm{mm}$ angibt. Ferner ist beim Okular zwar der von mir gemessene Abstand von Linse 1 und 4, nämlich $177{,}6\,\mathrm{mm}$, dem \textsc{Gauss}schen Werte $176{,}8\,\mathrm{mm}$ sehr benachbart, aber die Abstände der Linsen 1 und 2, 2 und 3, 3 und 4 von einander zeigen deutliche Differenzen. Statt der \textsc{Gauss}schen Werte
\[
h'=46{,}2\,\mathrm{mm},\quad h''=78{,}0\,\mathrm{mm},\quad h'''=52{,}6\,\mathrm{mm}
\]
\noindent finde ich
\[
h'=49{,}2\,\mathrm{mm},\quad h''=80{,}8\,\mathrm{mm},\quad h'''=47{,}6\,\mathrm{mm}.
\]
\noindent Es wäre allerdings nicht ausgeschlossen, daß an dem Okular nachträglich Verschiebungen der Linsen gegeneinander vorgenommen worden wären, so daß es nicht möglich ist, mit \textls[120]{Sicherheit} die Identität des
\par\HistoricalFoot{XI1.}{17}
\clearpage\SourceLeaf{135}\HistoricalHead{130}{L}{PHYSIK.\quad NACHLASS.}
\noindent vorliegenden Fernrohrs mit dem von \textsc{Gauss} in dem Abschnitt [II, B] untersuchten zu behaupten oder zu verneinen, zumal da aus einem Briefe von \textsc{Gauss} an \textsc{Encke} vom 2. Januar 1840 (siehe weiter unten S. 154) hervorzugehen scheint, daß \textsc{Gauss} mehrere \textsc{Fraunhofer}sche Instrumente von etwa gleicher Größe zur Verfügung gehabt hat.
\begin{center}\textls[120]{Zum Abschnitt [III.], S. 126.}\end{center}
\noindent\hspace*{\parindent}Zur Lösung der hier gestellten Aufgabe macht \textsc{Gauss} drei versuchsweise Ansätze, in der Handschrift als I\textsuperscript{te}, II\textsuperscript{te} und III\textsuperscript{te} Hypothese bezeichnet. Das Objektiv ist als achromatisch angenommen, die beiden Okularlinsen, wie dies auch heute noch zu geschehen pflegt, als aus demselben Glasflusse bestehend. Die drei Brennweiten sind gegeben. In den drei Hypothesen werden geeignet erscheinende Werte von $h$ und $h'$ angenommen, sodaß nach der Vorschrift von Abschnitt [I.] die Größen $\alpha,\beta,\gamma,\delta$ und aus diesen der farbige Rand bestimmt werden kann. Für die Berechnung des letzteren ist noch folgendes zu bemerken: Der farbige Rand ist nach Gleichung [VIIa] in Abschnitt [I.] proportional $\frac{\partial n}{n-1}$; da der Rand zum Verschwinden gebracht werden soll, so braucht der Proportionalitätsfaktor $\frac{\partial n}{n-1}$ gar nicht erst berechnet zu werden, und so hat \textsc{Gauss} ihn einfach weggelassen. Der Deutlichkeit halber ist er jedoch beim Abdruck in Klammern wieder hinzugesetzt worden.

Ein Vergleich des zweiten Ansatzes mit dem ersten zeigt, daß durch die zweite Wahl von $h$ und $h'$ der farbige Rand beträchtlicher geworden ist; bei der dritten Hypothese sind daher $h$ und $h'$ im umgekehrten Sinne geändert worden. In der Handschrift ist die Rechnung nicht durchgeführt; sie ist beim Abdruck in Klammern hinzugesetzt worden und zeigt, daß nun der farbige Rand rund zwanzigmal kleiner ist, als bei den vorhergehenden Annahmen, sodaß die dritte Hypothese als die sehr angenäherte Lösung der gestellten Aufgabe angesehen werden kann. Gleichzeitig überzeugt man sich, daß auch die für ein kurzsichtiges Auge von der Augenweite $k$ geltende Beziehung
\[
\beta''\delta''k=-1
\]
\noindent sehr angenähert erfüllt ist. — Die Angaben für den Ort des Auges, Blendung des Objektivs und Gesichtsfeld des Fernrohres rühren wieder von \textsc{Gauss} her und zeigen, daß er sich wohl selbst durch Rechnung überzeugt hatte, daß dieser dritte Ansatz dem gestellten Problem genügt.
\begin{flushright}\textsc{Schaefer.}\end{flushright}
\begin{center}\rule{.25\linewidth}{.4pt}\end{center}
\clearpage\endgroup
